\chapter{Modelo matemático implementado}

\section{Alcance e hipótesis}

Se implementó un modelo nominal computable con dos niveles. El modelo global
representó una base flotante de seis grados de libertad, doce coordenadas
articulares y fuerzas de contacto. El modelo reducido representó cada pata como
una cadena serial de tres grados de libertad para calcular cinemática,
jacobianos, dinámica local y trayectorias cartesianas. El vector general y el
vector de una pata se definieron como

\begin{equation}
\mathbf q=
\begin{bmatrix}\mathbf r_B^{\mathsf T}&\boldsymbol\eta_B^{\mathsf T}&
\mathbf q_J^{\mathsf T}\end{bmatrix}^{\mathsf T}\in\mathbb R^{18},
\qquad
\mathbf q_i=\begin{bmatrix}q_{c,i}&q_{f,i}&q_{t,i}\end{bmatrix}^{\mathsf T}.
\label{eq:model-levels}
\end{equation}

Se asumieron eslabones rígidos, suelo plano nominal, contacto unilateral,
servos controlados en posición y locomoción suficientemente lenta para analizar
el gateo con un criterio cuasiestático. La holgura, flexibilidad, fricción y
dispersión de fabricación no se consideraron parámetros identificados. Por
ello, la concordancia entre ecuaciones, URDF y simuladores verificó la
implementación, pero no validó físicamente el ejemplar.

Se utilizó la convención REP--103: $x$ hacia delante, $y$ hacia la izquierda y
$z$ hacia arriba. FL, FR, RL y RR identificaron las patas. El factor
$s_i=+1$ se empleó para las patas izquierdas y $s_i=-1$ para las derechas.

\section{Transformaciones y cinemática directa}

La transformación entre dos marcos se representó mediante

\begin{equation}
{}^A\mathbf T_B=
\begin{bmatrix}
{}^A\mathbf R_B & {}^A\mathbf p_B\\
0&0&0&1
\end{bmatrix}.
\label{eq:homogeneous-nova}
\end{equation}

Para una pata, $l_c$, $l_f$ y $l_t$ designaron las longitudes de coxa, fémur y
tibia. Sea $a=s_iq_c$ y $Q=q_f+q_t$. La posición del pie respecto de la cadera
se obtuvo como

\begin{align}
x &= -l_f\sin q_f-l_t\sin Q,\label{eq:nova-fkx}\\
z_p &= -l_f\cos q_f-l_t\cos Q,\label{eq:nova-fkzp}\\
y &= \cos(a)s_il_c-\sin(a)z_p,\label{eq:nova-fky}\\
z &= \sin(a)s_il_c+\cos(a)z_p.\label{eq:nova-fkz}
\end{align}

Estas expresiones se implementaron en \texttt{forward\_leg}. Para la postura
nominal $[q_c,q_f,q_t]=[0,10,0,42,-0,84]$ rad, con
$[l_c,l_f,l_t]=[0,090,0,105,0,132]$ m, la pata delantera izquierda produjo

\begin{equation}
{}^{H_{FL}}\mathbf p_{FL}=
\begin{bmatrix}0,0110095&0,1111545&-0,2063360\end{bmatrix}^{\mathsf T}\text{ m}.
\label{eq:nova-fk-example}
\end{equation}

Este punto fue la posición neutral usada para construir las trayectorias
cartesianas nominales. Su coincidencia con el código no convirtió las
longitudes nominales en mediciones físicas.

\section{Cinemática inversa y dominio}

Para una referencia cartesiana $\mathbf p=[x,y,z]^{\mathsf T}$ se eliminó
primero la rotación de coxa:

\begin{align}
z_p&=-\sqrt{y^2+z^2-l_c^2},\label{eq:nova-ikzp}\\
a&=\operatorname{atan2}(z,y)-\operatorname{atan2}(z_p,s_il_c),
&q_c&=s_ia.\label{eq:nova-ikcoxa}
\end{align}

Con $d=-z_p$, $b=-x$ y $r^2=b^2+d^2$, el triángulo fémur--tibia se resolvió
mediante

\begin{align}
D&=\frac{r^2-l_f^2-l_t^2}{2l_fl_t},&-1\leq D\leq1,\\
q_t&=-\arccos D,\label{eq:nova-knee}\\
q_f&=\operatorname{atan2}(b,d)-
\operatorname{atan2}\left(l_t\sin q_t,l_f+l_t\cos q_t\right).
\label{eq:nova-ikfemur}
\end{align}

La rama negativa de $q_t$ conservó la rodilla flexionada compatible con el
URDF. El controlador rechazó referencias sin solución real, fuera de límites o
discontinuas; no las saturó silenciosamente. La prueba de ida y vuelta FK--IK
recuperó la postura nominal con error numérico inferior a $10^{-12}$ rad.

\section{Jacobiano y singularidades}

La relación diferencial se definió mediante

\begin{equation}
\dot{\mathbf p}_i=\mathbf J_i(\mathbf q_i)\dot{\mathbf q}_i,
\qquad
\mathbf J_i=\frac{\partial\mathbf p_i}{\partial\mathbf q_i}.
\label{eq:nova-jacobian-definition}
\end{equation}

Con $A=l_f\sin q_f+l_t\sin Q$, $B=l_t\sin Q$ y $L=s_il_c$, se implementó

\begin{equation}
\mathbf J_i=
\begin{bmatrix}
0&z_p&-l_t\cos Q\\
s_i[-\sin(a)L-\cos(a)z_p]&-\sin(a)A&-\sin(a)B\\
s_i[\cos(a)L-\sin(a)z_p]&\cos(a)A&\cos(a)B
\end{bmatrix}.
\label{eq:nova-jacobian}
\end{equation}

El principio de potencia virtual produjo
$\boldsymbol\tau_i=\mathbf J_i^{\mathsf T}\mathbf f_i$. En la postura nominal
FL se obtuvieron valores singulares
$(0,252836,0,232233,0,038037)$ m/rad y número de condición $6,647$. La postura
no fue singular, aunque presentó una dirección de menor manipulabilidad. El
jacobiano analítico se contrastó con diferencias centrales y la conversión
fuerza--par conservó potencia virtual dentro de la tolerancia numérica.

\section{Dinámica nominal}

El modelo de base flotante se escribió como

\begin{equation}
\mathbf M(\mathbf q)\ddot{\mathbf q}+
\mathbf C(\mathbf q,\dot{\mathbf q})\dot{\mathbf q}+
\mathbf g(\mathbf q)+\boldsymbol\tau_f=
\mathbf S^{\mathsf T}\boldsymbol\tau+
\mathbf J_c^{\mathsf T}\boldsymbol\lambda.
\label{eq:nova-dynamics}
\end{equation}

La matriz $\mathbf S$ seleccionó las doce articulaciones actuadas. La matriz de
masa y el vector gravitacional se obtuvieron a partir de energía cinética y
potencial:

\begin{align}
K&=\frac12\sum_k\left[m_k\mathbf v_{C_k}^{\mathsf T}\mathbf v_{C_k}+
\boldsymbol\omega_k^{\mathsf T}\mathbf R_k\mathbf I_k
\mathbf R_k^{\mathsf T}\boldsymbol\omega_k\right],\\
U&=\sum_km_kgz_{C_k},\\
M_{ab}&=\frac{\partial^2K}{\partial\dot q_a\partial\dot q_b},&
g_a&=\frac{\partial U}{\partial q_a}.
\end{align}

Para una pata se implementó la dinámica inversa

\begin{equation}
\boldsymbol\tau_i=\mathbf M_i\ddot{\mathbf q}_i+
\mathbf c_i+\mathbf g_i+\boldsymbol\tau_{f,i}-
\mathbf J_i^{\mathsf T}\mathbf f_i.
\label{eq:nova-inverse-dynamics}
\end{equation}

En el ejemplo nominal, los autovalores de $\mathbf M_i$ fueron
$(0,00226642,0,00758520,0,00910534)$ kg m$^2$, todos positivos. Las pruebas
automatizadas comprobaron simetría, definición positiva, términos de Coriolis
nulos a velocidad cero y pares finitos. Estos resultados verificaron
propiedades internas del modelo, no la exactitud de los pares físicos.

\section{Actuador, contacto y estabilidad}

El MG996R se representó mediante una envolvente nominal lineal de par--velocidad:

\begin{equation}
|\tau_{\max}(\omega,V)|=\tau_{\mathrm{stall}}(V)
\max\left(0,1-\frac{|\omega|}{\omega_0(V)}\right).
\label{eq:nova-actuator}
\end{equation}

El par de bloqueo de catálogo no se interpretó como par continuo seguro. Cada
servo todavía requiere medición individual de centro, extremos, sentido,
corriente, velocidad, temperatura y respuesta bajo carga.

Para el contacto unilateral se emplearon las condiciones

\begin{align}
\phi_i&\geq0,&\lambda_{n,i}&\geq0,&\phi_i\lambda_{n,i}&=0,\\
\sqrt{\lambda_{x,i}^2+\lambda_{y,i}^2}&\leq\mu\lambda_{z,i}.
\label{eq:nova-contact}
\end{align}

Gazebo y MuJoCo resolvieron estas interacciones mediante sus propios algoritmos;
el modelo penalizado del proyecto no se presentó como reproducción exacta de
los solucionadores.

El centro de masa nominal se calculó como promedio ponderado de cuerpo y
eslabones. Para un polígono de contacto ordenado antihorariamente, la distancia
firmada a cada arista fue

\begin{equation}
d_j=\frac{e_{x,j}(c_y-p_{y,j})-e_{y,j}(c_x-p_{x,j})}
{\|\mathbf e_j\|},
\qquad m_s=\min_jd_j.
\label{eq:nova-static-margin}
\end{equation}

Un margen $m_s>0$ indicó que la proyección del centro de masa estaba dentro del
polígono. En postura nominal con cuatro apoyos se obtuvo
$m_s=0,073126$ m. Durante la marcha, el cálculo utilizó contactos medidos y se
declaró no disponible con menos de tres puntos no colineales.

\section{Trazabilidad y estado de validez}

La Tabla~\ref{tab:model-traceability} relaciona las familias matemáticas con su
implementación y comprobación.

\begin{table}[H]
\centering
\caption{Trazabilidad resumida del modelo matemático implementado.}
\label{tab:model-traceability}
\footnotesize
\begin{tabular}{>{\raggedright\arraybackslash}p{2.3cm}
>{\raggedright\arraybackslash}p{3.8cm}
>{\raggedright\arraybackslash}p{5.2cm}}
\toprule
Componente & Implementación & Evidencia de verificación \\
\midrule
FK e IK & \texttt{kinematics.py} & Ida y vuelta, dominio, continuidad y
alcanzabilidad de las marchas \\
Jacobiano & modelo matemático & Diferencias centrales y potencia
virtual \\
Dinámica & modelo matemático & Simetría, autovalores positivos,
Coriolis y pares finitos \\
Modelo físico & URDF/SDF y MJCF & Consistencia de geometría, masas, límites,
nombres y ejes \\
Contacto y margen & monitores de contacto y estabilidad & Sensores Gazebo,
pruebas unitarias y bolsas
de contacto \\
Marchas & generadores cartesianos de gateo y paso &
Alcanzabilidad, periodicidad, líneas base y repetición \\
\bottomrule
\end{tabular}
\end{table}

La suite completa alcanzó 50 pruebas aprobadas al incorporar el análisis de
contacto crudo y filtrado. El modelo detallado y sus ejemplos reproducibles se
conservaron además en \texttt{Documentacion/MODELO\_MATEMATICO\_LATEX}. El OE1
quedó avanzado en modelado y verificación computacional, pero permaneció
parcial hasta medir y contrastar el ejemplar físico completo.
